\uline{\textbf{Motivation initiale}}

Au départ, nous avions choisi le projet DISC-09, intitulé "Conception d'un dispositif IoT intelligent pour la sécurité des travailleurs isolés". Ce sujet nous a immédiatement attirés, car nous cherchions un défi technique réunissant informatique et systèmes embarqués, tout en répondant à une problématique humaine forte. L'objectif était de créer un outil pour sécuriser les ouvriers travaillant dans des conditions extrêmes ou isolées, en fusionnant des données de capteurs pour détecter des risques critiques comme les chutes ou les gaz toxiques. 

\bigskip

\uline{\textbf{Prise de contact avec des professionnels}}

Nous avons pris contact avec Monsieur Luc-Géry Helle, directeur régional génie civil de GTM Ouest chez VINCI Construction France, afin de parler de notre projet et de poser des questions pour se renseigner sur le domaine du bâtiment et des travaux publics. Nous souhaitions bénéficier de conseils de professionnels dans le but de rendre notre projet le plus réaliste possible. 
Monsieur Helle nous a proposé une visioconférence pour en discuter. Nous lui avons envoyé les documents propres à notre projet. Lors de cette réunion, il nous informe d'une contrainte réglementaire : la législation française interdit ou encadre de manière extrêmement stricte le travail isolé sur les chantiers de construction. Les procédures d'entreprise imposent une co-présence systématique pour des raisons de sécurité. Notre projet s'avère donc moins pertinent. Nous souhaitons modifier notre projet afin qu'il contribue concrètement à un enjeu de sécurité sur les chantiers de construction. Monsieur Helle nous apprend que la première cause d'accidents mortels sur les chantiers sont les collisions entre les engins de chantier et les ouvriers piétons. En effet, il y en a environ quinze morts chaque année sur les chantiers selon l'OPPBTP. Les engins étant larges et imposants, le conducteur, concentré sur sa tâche, peut ne pas parfaitement avoir conscience de l'environnement qui l'entoure. Il nous apprend qu'actuellement, il n'y a pas encore de solution efficace en toutes circonstances. Ils investissent dans des caméras pour leurs engins. Pour appuyer l'importance du sujet, Monsieur Helle nous annonce que la Safety Week 2026 organisée par Vinci Construction porte sur le sujet de la collision engin-piéton. 

\bigskip

\uline{\textbf{Analyse des contraintes}}

Suite à ces échanges, nous avons repensé notre projet pour qu'il réponde au mieux aux vrais besoins de l'entreprise. Nous imaginons alors un système de capteurs à mettre sur les engins de chantier pour détecter les ouvriers et prévenir en cas de danger de collisions. Un écran présent dans la cabine de l'engin permettrait au conducteur de localiser en direct la position des ouvriers. Certaines technologies impliquent qu'un capteur soit dans un boîtier porté obligatoirement par chaque ouvrier, d'autres n'en nécessitent pas. 
Monsieur Helle nous indique que si chaque ouvrier doit porter un boîtier, cela serait compliqué à grande échelle de former les employés, de vérifier qu'ils le portent correctement tout au long de leur journée ainsi que si les boîtiers sont bien chargés. Dans la réalité, les employés se dépêchent une fois leur journée terminée et pourraient oublier de mettre à charger leur boîtier. Cependant, Monsieur Helle nous encourage tout de même, car cette option reste intéressante si elle fonctionne bien. Il nous conseille cependant de ne pas intégrer de bouton pour activer ou désactiver ce boîtier. En effet, le bouton risque de s'endommager rapidement et de compromettre l'étanchéité du dispositif.
D'autres conseils nous sont donnés. Il n'est pas nécessaire que le système fonctionne quand l'engin est à l'arrêt, un simple bouton dans la cabine du chauffeur pour allumer et éteindre le système suffit. Il faut que le système soit hermétique et résistant aux conditions difficiles sur chantiers : mauvais temps, eau, boue, poussière, écrasements, chocs. L'écran dans la cabine de l'engin n'est pas jugé obligatoire, une alarme sonore suffit pour prévenir d'un danger. Cependant, la distance de sécurité minimale varie en fonction de l'engin alors cela pourrait être intéressant à pouvoir la définir à souhait en fonction de l'engin. Dans ce cas, il faut s'assurer que l'écran soit résistant et développer une interface web permettant de personnaliser cette zone de sécurité.
L'entreprise souhaite un produit le plus personnalisable possible. Il faut tout de même faire attention au risque de modification irresponsable qui pourrait mettre en danger les travailleurs tout comme le matériel et l'entreprise.

\bigskip

\uline{\textbf{Surveillance des travailleurs}}

Le sujet sensible de la surveillance au travail a aussi été abordé avec le directeur régional de GTM Ouest et traité dans notre rapport éthique. Nous avions initialement imaginé développer une interface web où la position exacte de chaque ouvrier serait consultable en direct sur une carte dans le but de localiser les potentiels accidents. Cependant, cela pourrait être une forme d'atteinte à la vie privée des travailleurs, qui se sentiraient en permanence observés. Nous avons donc abandonné cette idée. A la place, nous voulons mettre en place un historique des détections de pénétration dans la zone de sécurité, stocké en local, avec l'heure, la date ainsi que l'identifiant du boîtier.

\bigskip

\uline{\textbf{Visite de chantier}}

Monsieur Helle nous a proposé de nous rendre sur un chantier afin de nous rendre compte des conditions réelles dans lesquelles travaillent les ouvriers. Il nous a mis en contact avec Monsieur Mathieu Baré, directeur prévention délégation génie civil Ouest chez VINCI Construction France. Seules trois personnes pouvaient se rendre sur le chantier : Maxime, Clarisse et Benoît. Le 2 juin 2026, accompagné du directeur prévention, nous nous sommes rendus sur le chantier de la gare de Nanterre La Folie de la future ligne de métro 15 du Grand Paris Express.
Nous avons pu entrer dans les bases vie, lieux où les ouvriers se préparent en arrivant sur le chantier et prennent leurs pauses. Les équipements de protection individuelle sont stockés dans les bases vie et sont obligatoires sur le chantier. Nous avons aussi été équipés d'EPI : chaussures de sécurité, gants, gilets fluorescents de sécurité, lunette de protection, bouchon d'oreilles et casque de chantier. Nous avons ainsi remarqué qu'il est possible d'intégrer un boîtier sur le casque des ouvriers.
Nous sommes descendus sur le chantier au niveau du tunnel. Nous avons pu nous approcher au plus près d'engins équipés de certains dispositifs afin d'éviter les collisions engins-piétons. Ils nous ont présenté un engin équipé de laser permettant de délimiter une zone de sécurité à même le sol autour de l'engin. Cette solution est efficace dans un environnement sombre, mais devient sans effet lorsque l'environnement est bien lumineux. Une pelleteuse est équipée d'une caméra avec intelligence artificielle dont l'image est diffusée sur un écran dans la cabine de l'engin. Un système de zone de danger modulaire est utilisé pour varier les types d'alerte à envoyer selon la distance. Le conducteur peut alors en bénéficier pour s'assurer de l'environnement qui l'entoure afin d'éviter un accident. Cependant, la caméra ne permet pas toujours de voir chaque angle mort. De plus, la vue 360° déforme notre perception des distances. 
Ainsi, Monsieur Baré insiste sur le fait que plusieurs dispositifs sont mis en place pour éviter les accidents entre engins et piétons, mais chacun présente des avantages et des inconvénients et sont plus ou moins efficaces selon les conditions et situations. Aucun système à ce jour ne permet de supprimer complètement ce risque.
